% GENERATED FILE. Edit canonical lessons, then run npm run generate:books.
% canonical-source-hash: fnv1a64:e10ece4382823bc3
% canonical-lessons: LA-C287-laetus-eram, LA-C287-spero, LA-R287-how-it-felt

\chapter{How It Felt}
\label{ch:la-how-it-felt}
\hlchaptermodality{287}

\begin{chapteropening}
\textbf{By the end of this chapter:} \emph{I can say how I felt, with laetus eram, and what I hope for and long to do.}

Complete LA-R287-how-it-felt: I can say how I felt, with laetus eram, and what I hope for and long to do.
\end{chapteropening}

\section[laetus eram, laeta eram]{laetus eram — how it felt}

\label{lesson:LA-C287-laetus-eram}



\begin{quote}
Say \textquotedblleft{}happy\textquotedblright{}, then \textquotedblleft{}sad\textquotedblright{}. (\emph{Laetus. Trīstis.}) Then \textquotedblleft{}I was\textquotedblright{}, the open past. (\emph{Eram.})
\end{quote}

\begin{grammarlens}[title={Grammar Lens: the feeling agrees with who felt it}]
A feeling in a story is a state, so it takes the imperfect, and the adjective matches the person who felt it:

\begin{quote}
\textbf{Laetus eram.} — I was happy (a man speaking). \textbf{Laeta eram.} — I was happy (a woman speaking). \textbf{Sollicitī erāmus.} — We were worried. \textbf{Trīstis erat.} — He was sad, or she was: \emph{trīstis} serves both.
\end{quote}

\emph{Anxius}, \emph{sollicitus} and \emph{īrātus} follow \emph{laetus}: \emph{-us} for a man, \emph{-a} for a woman, \emph{-ī} for several. \emph{Trīstis} has one form for a man and a woman, and \emph{trīstēs} for several.
\end{grammarlens}

\begin{culture}
The cause of a feeling takes the \emph{quod} you met as a reader, and here both verbs are scene:

\begin{quote}
\textbf{Laeta eram, quod caelum serēnum erat.} — I was happy, because the sky was clear. \textbf{Trīstis eram, quod pluēbat.} — I was sad, because it was raining.
\end{quote}

Latin can also put the feeling in a verb: \emph{gaudēbam}, \textquotedblleft{}I was glad\textquotedblright{}, beside \emph{laetus eram}; \emph{timēbam}, \textquotedblleft{}I was afraid\textquotedblright{}. The adjective describes the person; the verb tells what the person was doing inside.
\end{culture}

\begin{cousinweb}
Two of these feelings are bodily, and their own lessons said how. \textbf{anxius} is built on \emph{angō}, \textquotedblleft{}to squeeze, to choke\textquotedblright{} (Lewis \& Short): worry as a tight throat. \textbf{sollicitus} is built from \emph{sollus}, \textquotedblleft{}whole\textquotedblright{}, and \emph{cieō}, \textquotedblleft{}to stir up\textquotedblright{} (Lewis \& Short): stirred through and through. English took them over as \emph{anxious} and \emph{solicitous}.
\end{cousinweb}

\subsection*{Your turn}
Say these aloud:

\begin{itemize}
\raggedright
  \item \emph{Laetus eram.} — then a woman says it: \emph{Laeta eram.}
  \item \emph{Sollicitī erāmus.}
  \item \emph{Laeta eram, quod caelum serēnum erat.}
  \item \emph{Trīstis eram, quod nāvis nōn veniēbat.} — I was sad, because the ship was not coming
  \item \emph{In nāve īrātus erat, et timēbam.}
  \item how you felt yesterday, and why
\end{itemize}

\begin{tcolorbox}[breakable,colback=teal!4,colframe=teal!35!black,title={Before you move on}]
A woman says \textquotedblleft{}I was happy\textquotedblright{}. (\emph{\textbf{Laeta eram.}}) \textquotedblleft{}We were worried\textquotedblright{}? (\emph{\textbf{Sollicitī erāmus.}}) Which feeling word has one form for a man and a woman? (\emph{\textbf{Trīstis.}}) Which tense does a feeling in a story take? (The \textbf{imperfect}: it is a state.)
\end{tcolorbox}
\section[spērō, spērāre]{spērō — I hope}

\label{lesson:LA-C287-spero}



\begin{quote}
Say \textquotedblleft{}I was happy\textquotedblright{}. (\emph{Laetus eram.}) Then the noun for \textquotedblleft{}hope\textquotedblright{}. (\emph{Spēs.})

A feeling looks back. A hope looks forward, and Latin has a verb for it built on the noun you own.
\end{quote}

\subsection*{You'll want to know: spērō}
\textbf{spērō, spērāre} — \textquotedblleft{}to hope\textquotedblright{}. It takes three shapes.

\begin{quote}
\textbf{Pācem spērō.} — I hope for peace. \textbf{Ut spērō, valēs.} — As I hope, you are well. \textbf{Spērō tē valēre.} — I hope you are well.
\end{quote}

The first hopes for a thing, in the accusative. The second tucks \emph{ut spērō}, \textquotedblleft{}as I hope\textquotedblright{}, into a sentence; Cicero writes it in his letters. The third is the accusative and infinitive of \emph{putō}: \emph{tē}, then \emph{valēre}.

When the hope is about later, Latin prefers an infinitive made for the future, a form still to come. For now, hope for things, and for how things are.

\begin{grammarlens}[title={Grammar Lens: cupiō, volō + infinitive}]
What you long \textbf{to do} goes with \textbf{cupiō}, \textquotedblleft{}I long\textquotedblright{}, or \textbf{volō}, \textquotedblleft{}I want\textquotedblright{}, and a plain infinitive:

\begin{quote}
\textbf{Rōmam vidēre cupiō.} — I long to see Rome. \textbf{Mercātor esse volō.} — I want to be a merchant.
\end{quote}

\emph{Mercātor} stays nominative: it is the same person as the one who wants. And \textbf{optō}, \textquotedblleft{}to wish for, to choose\textquotedblright{}, hopes for a thing, like \emph{spērō}: \emph{nāvem optō}.
\end{grammarlens}

\begin{cousinweb}
\textbf{spērō} is the verb of \emph{spēs}, \textquotedblleft{}hope\textquotedblright{}, and EtymDB 2.1 traces it, like the noun, to Proto-Indo-European \emph{*speh\textsubscript{1}-}, \textquotedblleft{}to prosper, to turn out well\textquotedblright{}. Put \emph{dē-}, \textquotedblleft{}away\textquotedblright{}, in front and you have \emph{dēspērō}, \textquotedblleft{}to give up hope\textquotedblright{} (Lewis \& Short), which reached English as \emph{despair} and \emph{desperate}.
\end{cousinweb}

\subsection*{Your turn}
Say these aloud:

\begin{itemize}
\raggedright
  \item \emph{spērō, spērāre}
  \item \emph{Pācem spērō.} — then \emph{Spērō tē valēre.}
  \item \emph{Ut spērō, valēs.}
  \item \emph{Rōmam vidēre cupiō.}
  \item \emph{Nauta esse volō, et nāvem optō.}
  \item one thing you long to do, with \emph{cupiō}
\end{itemize}

\begin{tcolorbox}[breakable,colback=teal!4,colframe=teal!35!black,title={Before you move on}]
Say \textquotedblleft{}I hope you are well\textquotedblright{}. (\emph{\textbf{Spērō tē valēre.}}) Say \textquotedblleft{}I long to see Rome\textquotedblright{}. (\emph{\textbf{Rōmam vidēre cupiō.}}) Which English word grew from this verb's family? (\textbf{Despair.})
\end{tcolorbox}
\section[(review)]{Review — how it felt, and what you hope for}

\label{lesson:LA-R287-how-it-felt}



\begin{quote}
A woman says \textquotedblleft{}I was worried\textquotedblright{}. (\emph{Sollicita eram.}) Then say \textquotedblleft{}I hope for peace\textquotedblright{}. (\emph{Pācem spērō.})
\end{quote}

\begin{grammarlens}[title={What you've built}]
Two lessons, and an experience has a before and an after. Before: a feeling with \emph{eram}, agreeing with the one who felt it, and its cause with \emph{quod}. After: \emph{spērō} for what you hope for, and \emph{cupiō} or \emph{volō} with an infinitive for what you long to do.

\begin{quote}
\emph{Sarcinam āmīsī, et trīstis eram. Sed nunc laetus sum, et Rōmam vidēre cupiō.} — I lost my pack, and I was sad. But now I am happy, and I long to see Rome.
\end{quote}

An event, a feeling in the past, a feeling now, and a wish.
\end{grammarlens}

\subsection*{Your turn}
Say these aloud:

\begin{itemize}
\raggedright
  \item \emph{Laeta eram, quod caelum serēnum erat.}
  \item \emph{Quid accidit? — Nāvis discessit, et sollicitī erāmus.}
  \item \emph{Spērō tē valēre.}
  \item \emph{Sarcinam āmīsī, et trīstis eram. Sed nunc Rōmam vidēre cupiō.}
  \item a day that went wrong, how you felt, and what you hope for now
\end{itemize}

\begin{tcolorbox}[breakable,colback=teal!4,colframe=teal!35!black,title={Before you move on}]
Say \textquotedblleft{}we were worried\textquotedblright{}. (\emph{\textbf{Sollicitī erāmus.}}) Say \textquotedblleft{}I hope you are well\textquotedblright{}. (\emph{\textbf{Spērō tē valēre.}}) Say \textquotedblleft{}I long to see Rome\textquotedblright{}. (\emph{\textbf{Rōmam vidēre cupiō.}}) Say \textquotedblleft{}I have lost my pack\textquotedblright{}. (\emph{\textbf{Sarcinam āmīsī.}})
\end{tcolorbox}
